%%
% Template for Assignment Reports
% 
%

\documentclass[11pt]{article}

\usepackage{fancyhdr} % Required for custom headers
\usepackage{lastpage} % Required to determine the last page for the footer
\usepackage{extramarks} % Required for headers and footers
\usepackage{graphicx,color}
\usepackage{anysize}
\usepackage{amsmath}
\usepackage{natbib}
\usepackage{caption}
\usepackage{hyperref}
\usepackage{listings}
\usepackage{float}
%\usepackage[strings]{underscore}
\usepackage{url}

% Margins
%\topmargin=-0.45in
%\evensidemargin=0in
%\oddsidemargin=0in
\textwidth=6.5in
%\textheight=9.0in
%\headsep=0.25in 
%\linespread{1.0} % Line spacing

%%------------------------------------------------
%% Image and Listing code
%%------------------------------------------------
%% Examples for the commands in the document below
%%
%% includecode:
%% \includecode{caption for table of listings}{caption for reader}{filename}
%% - includes a file with code and adds a caption that should describe the code in some detail and a shorter caption for the table of listings
\newcommand{\includecode}[4]{\lstinputlisting[float,floatplacement=H, caption={[#1]#2}, captionpos=b, frame=single, label={#3}]{#4}}


%% includescalefigure:
%% \includescalefigure{label}{short caption}{long caption}{scale}{filename}
%% - includes a figure with a given label, a short caption for the table of contents and a longer caption that describes the figure in some detail and a scale factor 'scale'
\newcommand{\includescalefigure}[5]{
\begin{figure}[htb]
\centering
\includegraphics[width=#4\linewidth]{#5}
\captionsetup{width=.8\linewidth} 
\caption[#2]{#3}
\label{#1}
\end{figure}
}

%% includefigure:
%% \includefigure{label}{short caption}{long caption}{filename}
%% - includes a figure with a given label, a short caption for the table of contents and a longer caption that describes the figure in some detail
\newcommand{\includefigure}[4]{
\begin{figure}[htb]
\centering
\includegraphics{#4}
\captionsetup{width=.8\linewidth} 
\caption[#2]{#3}
\label{#1}
\end{figure}
}


%%------------------------------------------------
%% Parameters
%%------------------------------------------------
% Set up the header and footer
\pagestyle{fancy}
\lhead{\authorName} % Top left header
\rhead{\moduleCode\ - \shortAssignmentTitle} % Top center header
%\rhead{\firstxmark} % Top right header
\lfoot{\lastxmark} % Bottom left footer
\cfoot{} % Bottom center footer
\rfoot{Page\ \thepage\ of\ \pageref{LastPage}} % Bottom right footer
\renewcommand\headrulewidth{0.4pt} % Size of the header rule
\renewcommand\footrulewidth{0.4pt} % Size of the footer rule

\setlength\parindent{0pt} % Removes all indentation from paragraphs
\newcommand{\assignmentTitle}{Assignment\ \#3: Biography of an influential software engineer} % Assignment title
\newcommand{\shortAssignmentTitle}{Assignment\ \#3: Biography} % Assignment title
\newcommand{\moduleCode}{CS3012} 
\newcommand{\moduleName}{Software\ Engineering} 
\newcommand{\authorName}{Jevgenijus\ Cistiakovas} % Your name ***EDIT HERE***
\newcommand{\authorID}{} % Your student ID ***EDIT HERE***
\newcommand{\reportDate}{October 16, 2019}


%%------------------------------------------------
%%	Title Page
%%------------------------------------------------
\title{
\vspace{-1in}
\begin{figure}[!ht]
\flushleft
\includegraphics[width=0.4\linewidth]{reduced-trinity.png}
\end{figure}
\vspace{-0.5cm}
\hrulefill \\
\vspace{0.5cm}
\textmd{\textbf{\moduleCode\ \moduleName}}\\
\textmd{\textbf{\assignmentTitle}}\\
\vspace{0.5cm}
\hrulefill \\
}
\author{\textbf{\authorName,\ \authorID}}
\date{\reportDate}


%%------------------------------------------------
%% Document
%%------------------------------------------------
\begin{document}
%% Defaults for listings
\lstset{language=Java, captionpos=b, frame=single}
\captionsetup{width=.8\linewidth} 

\maketitle
\tableofcontents
\vspace{0.5in}

%% We will skip a couple of components of reports such as abstracts, literature review, etc for the reports on assignments.
%%------------------------------------------------
\section{Introduction}
\label{sec:Intro}
This document is a biographical essay about a key software engineer. In particular, the subject of this biography is \textbf{Howard G. "Ward" Cunningham} - an American software engineer who developed the first wiki - a Web-based  software  application  that  allows  users  to  collaboratively  contribute  and  edit  articles  on  various  topics.\cite{cs-encyclopedia} He also pioneered in design patterns and extensively contributed to the field of programming methods, including the use of design patterns that became known as "extreme programming".
 
\section{Howard G. "Ward" Cunningham - Biography}
\subsection{Early life}
\par
Born on may 26, 1949, in Indiana, Cunningham grew up in an era before the Internet. He was however fully engaged with amateur radio - the next-best communications medium of that era. In particular, he was fascinated by the creativity of global community that gathered around the amateur radio. Cunningham learnt to program in high school, which had a special program to provide students access to mainframe computers at the Illinois Institute of Technology.\cite{wiki-revolution}\cite{innovators} He then attended Purdue University, where he received a bachelor’s degree in Electrical Engineering and Computer  Science,  and  then  a  master’s  in  Computer  Science. After  graduation  Cunningham  worked  as  a  researcher  in microcomputer  systems  for  Tektronix, a company best known for manufacturing measurement devices such as oscilloscopes. \cite{cs-encyclopedia}
\subsection{Wiki}
\par 
At Tektronix, Cunningham was assigned to keep track of projects, a task similar to what Berners-Lee faced when he went to CERN. To do this he modified a software product developed by  Apple. It was called HyperCard, and it allowed users to make their own hyperlinked cards and documents on their computers. Apple had little idea what to do with the software, so at Bill Atkinson’s insistence Apple gave it away free with its computers.\cite{innovators}\cite{c2-wiki}
\par
Using HyperCard, Cunningham built an application that allowed users to add free-form data to a database and link it to other entries by clicking a button. Users who tried it were fascinated by its potential. Cunningham then wanted to expand  it so users  could  access  it  over  networks, and one colleague suggested using the World Wide Web.\cite{cs-encyclopedia} He then created an Internet version of his HyperText program, writing it in just a few hundred lines of Perl code. The result was a new content management application that allowed users to edit and contribute to a Web page. Cunningham used the application to build a service, called the Portland Pattern Repository, that allowed software developers to exchange programming ideas and improve on the patterns that others had posted.\cite{innovators} As any tool, it needed a name. At first he thought of calling it QuickWeb. He  then  remembered  hearing  the  phrase wiki wiki or “quickly, quickly” in Hawaii, and he decided to call his system wikiwikiWeb. Today, it is just known as a wiki.\cite{cs-encyclopedia}\cite{c2-wiki}
\par 
WardsWiki (as it became known) allowed anyone to edit and contribute, without even needing a password. Previous versions of each page would be stored, in case someone spoiled one up, and there would be a “Recent Changes” page so that Cunningham and others could keep track of the edits. But there would be no supervisor or admin pre-approving the changes. Cunningham was sure that it would work, because he believed that people are generally good. It was just what Berners-Lee had envisioned, a Web that was read-write rather than read-only. “Wikis were one of the things that allowed collaboration,” Berners-Lee said. “Blogs were another.”\cite{innovators}
\subsection{Wikipedia}
\par
Like Berners-Lee, Cunningham made his basic software available for anyone to modify and use. Consequently, there were soon large number of wiki sites as well as open-source improvements to his software. But the wiki concept was not widely known beyond software engineers until January 2001, when Wikipedia was launched by Kimmy Wales and Larry Sanger.\cite{innovators}\cite{wiki-wikipedia}

Wikipedia became to Web content what GNU/Linux was to software: a peer-to-peer commons collaboratively created and maintained by volunteers who worked for the satisfactions they found in helping the community. It was a delightful, counter-intuitive concept, perfectly suited to the philosophy, attitude, and technology of the Internet. Anyone could edit a page, and the results would show up instantly.\cite{innovators}

\subsection{Design Patterns}
Ward Cunningham has contributed to the practice of object-oriented programming, in particular the use of pattern languages and (with Kent Beck) the class-responsibility-collaboration(CRC) cards. He also contributes to the extreme programming software development methodology. Much of this work was done collaboratively on the first wiki site. Cunningham is also an important figure in Agile movement. As a key figure in extreme programming he was one of the seventeen reputable software engineers that came up with a \textit{Manifesto for Agile Software Development}. The Manifesto played an important role in popularizing the Agile approach to software development. \cite{agile-manifesto}\cite{wiki-cunningham}
\subsection{Open Source}
\par
Cunningham was and is an important figure in Open Source. He never patented his main creation - wiki. Instead he made the wiki software open and available to everyone. He was happy to see  Wales and Sanger adapting his software for such an ambitious project as Wikipedia.\cite{innovators}
Cunningham worked for a few years on open-source projects at Microsoft. Later he held the position of a Director of Committer Community Development at the Eclipse Foundation - a non-profit organisation that oversees development of an open-source development environment - Eclipse.
\subsection{Present}
As Cunningham  \& Cunningham  Inc., Ward and  his  wife consult people in the area of object oriented programming.

%%------------------------------------------------------------------------------------

\section{Discussion}
Ward Cunningham is undoubtedly an influential figure in software engineering, who is unfortunately not well known to the people outside of software engineering. Despite having being a user of wikis and Eclipse for long time and knowing about extreme programming and design patterns, I only learnt about Ward Cunningham as one of the authors of the Agile Manifesto. The fact that I have used, learnt about and practised his inventions and ideas was a direct indicator to me of how influential person he is.
\par
The impact of the products on which Cunningham had worked is enormous. According to multiple surveys\cite{hp-agile-report}\cite{pmi-agile-report}\cite{state-of-agile-report}, most of the companies nowadays that develop software, as well as those that do not, make use of at least some practices of Agile development, including the extreme programming implementation of agile values. This includes but is not limited to short development cycles, extensive code reviews, unit testing of all code, etc.
With respect to design  patterns, Cunningham was the chairperson of the first ever conference on pattern languages at the University of Illinois at Urbana Champaign in 1994.\cite{cs-encyclopedia} The purpose of this conference is to develop and refine the art of software design patterns.
The biggest product of Eclipse Foundation - Eclipse IDE is one of the two most widely used IDEs for Java with proprietary IntelliJ IDEA being its main competitor.\cite{rebellabse-report} Similarly, the biggest implementation of wiki concept - Wikipedia is ranked as the 9th most visited site on the web by Alexa rankings. Wikipedia has hundred thousands active editors\cite{wiki-community}, with almost all of them being volunteers, exactly as Ward Cunningham has envisioned. Finally, the work of Cunningham, being it wiki or agile practices, always motivated people to collaborate in order to achieve the greater result.
\par
What fascinates me in Ward Cunningham is that unlike many of software engineers he was never driven by financial gain. His main creation - wiki - is free and open for all. Eclipse Foundation is a non-profit corporation. Design patterns, extreme programming and agile are generic practices intended for anyone to use. I think that there is an explanation to that. It looks to me like Cunningham's interest to open source is due to his involvement in amateur radio back in his childhood. When making wiki being a platform that is free and where everyone can contribute, he could have been recalling the global amateur radio community.
\par
The story of Ward Cunningham is also an example of importance of introducing children to technology at an early age. Cunningham interest in computers and software was conceived back in high school, when his teacher allowed him and other children to play with a mainframe via Teletype.

%%I find it interesting that Apple's HyperCard had influenced the  

%%Finally, Ward Cunningham's ideas of wiki and CRC cards could be traced back to HyperCard stack he wrote at work.\cite{c2-wiki} HyperCard also influenced the development of the Web. 
\par


I strongly believe that such humble and hard-working people like Cunningham should be examples for generations of new software engineers.

%% -Wiki
%% -Agile manifesto
%% -Extreme programming
%% -Eclipse foundation
%% -Design patterns

\section{Sources}
Most of the information about the life of Ward Cunningham and history of wiki is taken from \cite{innovators} \cite{wiki-revolution} \cite{cs-encyclopedia} as well as from the Wikipedia website.

%\bibliographystyle{apalike2}
\bibliographystyle{plain}
\bibliography{sources} 

\end{document}

